\documentclass[11pt]{article}

\usepackage[margin=1in]{geometry}
\usepackage{array}
\usepackage{longtable}
\usepackage{hyperref}

\title{Cachepawl Path C: Planner-Level Advisory Diagnosis of Hybrid Attention/Mamba Cache Overestimation}
\author{Formatting TODO: add venue-specific author and anonymization metadata}
\date{Venue-neutral draft}

\begin{document}
\sloppy

\maketitle

\begin{abstract}
Hybrid Attention/Mamba language models stress inference cache planners because
attention KV pages and Mamba state caches have different shapes and reuse
contracts. Cachepawl Path C evaluates whether these planner-level
inefficiencies can be diagnosed from observed vLLM cache-plan artifacts without
modifying the serving stack.

We present an observe/advisory workflow that captures vanilla
\texttt{vllm==0.21.0} cache artifacts for
\texttt{Zyphra/Zamba2-2.7B-instruct}, translates them into a Cachepawl schema,
replays the vLLM planner stage on real planner inputs, and emits an advisory
report through \texttt{cachepawl diagnose-vllm}. The planner-stage replay
matched the runtime scheduler cache configuration and did not change the
runtime scheduler state. Across a bounded 4-cell matrix for one model
(\texttt{max\_model\_len} in \{2048, 4096\},
\texttt{gpu\_memory\_utilization} in \{0.6, 0.7\}, and
\texttt{max\_num\_seqs=1}), estimated advisory savings ranged from
685,011,456 to 1,347,563,520 bytes. The overestimation ratio stayed
1.7333734577189286 and the wasted fraction stayed 0.4230902777777778 across
completed cells.

We also include a separate RTX 3060 planner-only pack for synthetic hybrid
KV/State workloads. In that deterministic pack, \texttt{cachepawl-avmp} stays
closer to useful bytes than the \texttt{vllm-style-padded} planner baseline,
but the result is not a vLLM serving measurement.

This is a diagnostic/advisory systems artifact. It does not claim runtime cache
substitution, allocator replacement, serving-time VRAM reduction, throughput
improvement, latency improvement, accuracy improvement, or quality impact.
\end{abstract}

\section{Introduction}

Hybrid Attention/Mamba language models place two different cache shapes behind
one inference runtime. Attention layers use KV blocks whose economics are tied
to token sequence growth, while Mamba or SSM layers use state tensors whose
useful byte footprint can diverge from the page shape used by a uniform cache
planner. Prior serving systems such as vLLM's PagedAttention
\cite{kwon2023pagedattention} and SGLang \cite{zheng2024sglang} provide the
systems context for this problem, while hybrid model families such as Jamba
\cite{lieber2024jamba} and Zamba2 \cite{glorioso2024zamba2} motivate the mixed
cache-shape setting. A cache planner that reserves by a shared page/block
abstraction can therefore allocate capacity according to a shape that is larger
than the useful state payload for hybrid layouts.

Cachepawl Path C studies this gap as an observe/advisory problem. The goal of
this phase is not to replace vLLM's allocator, rewrite live runtime state, or
report serving results. Instead, the goal is to answer a narrower systems
question: can Cachepawl observe a vanilla vLLM hybrid cache plan, translate it
into a stable schema, replay the planner stage on the observed inputs, and emit
a practical diagnostic report that quantifies planner-level overestimation?

\subsection{Scope of This Version}

This version is scoped to evidence-backed diagnosis. It includes a bounded Path
C matrix for one observed \texttt{Zyphra/Zamba2-2.7B-instruct} setup, a
separate deterministic RTX 3060 planner-only comparison for synthetic
\texttt{jamba-1.5-mini} workloads, and the \texttt{cachepawl diagnose-vllm}
artifact-input CLI. It excludes runtime allocator replacement, runtime cache
substitution, live VRAM reduction, serving throughput, serving latency,
quality, accuracy, and controlled substitution readiness.

The answer from the current artifact set is yes, within a bounded scope. The
evaluation uses \texttt{Zyphra/Zamba2-2.7B-instruct} under vanilla
\texttt{vllm==0.21.0} and a 4-cell advisory matrix with
\texttt{max\_model\_len} in \{2048, 4096\},
\texttt{gpu\_memory\_utilization} in \{0.6, 0.7\}, and
\texttt{max\_num\_seqs=1}. Across all four completed cells, the planner-level
overestimation ratio stayed 1.7333734577189286 and the planner-level wasted
fraction stayed 0.4230902777777778. Estimated advisory savings ranged from
685,011,456 to 1,347,563,520 bytes.

Those values are advisory metrics over planner artifacts. They are not runtime
VRAM measurements, throughput measurements, latency measurements, accuracy
measurements, or quality measurements. The phase deliberately keeps
artifact-input diagnosis separate from controlled substitution.

\subsection{Contributions}

The contributions of this report are evidence-bounded:

\begin{itemize}
\item a bounded method for capturing and translating vLLM hybrid cache-plan
artifacts;
\item planner-stage replay evidence that matches the runtime scheduler config;
\item a 4-cell Path C advisory matrix showing planner-level over-reservation
for one \path{Zyphra/Zamba2-2.7B-instruct} setup;
\item a deterministic RTX 3060 planner-only comparison showing
\texttt{cachepawl-avmp} stays closer to useful bytes than a uniform padded
planner baseline on three synthetic hybrid workloads;
\item a user-facing \texttt{cachepawl diagnose-vllm} artifact-input report that
is read-only and non-mutating;
\item an explicit contract checklist that keeps controlled substitution in
future work until Mamba state-index and state tensor paths are observable or
supported upstream.
\end{itemize}

\section{Method}

The Path C method is an observe/translate/replay/diagnose workflow. It is
designed to preserve the boundary between advisory analysis and runtime
mutation.

\subsection{Observe Vanilla vLLM Artifacts}

The workflow starts from a vanilla \texttt{vllm==0.21.0} run for
\texttt{Zyphra/Zamba2-2.7B-instruct}. Observation scripts collect cache-plan
artifacts and safe runtime metadata. The observation path is read-only: it does
not edit vLLM source, monkeypatch vLLM internals, replace allocators, or return
Cachepawl plans to vLLM.

The captured artifacts include the runtime cache plan, safe metadata about the
planner inputs, scheduler/KV manager structure, worker tensor layout, live
request/block assignment, and attention-side block-table/view metadata.

\subsection{Translate Cache Plans}

The runtime cache plan is translated into a Cachepawl schema that separates the
fields needed for advisory metrics from the fields that would be required for
future controlled substitution. The translation records cache group count,
cache tensor count, layer count, block count, per-group cache kind, page size,
block size, useful bytes, and observed reserved bytes.

\subsection{Replay the Planner Stage}

The workflow replays \texttt{vllm.v1.core.kv\_cache\_utils.get\_kv\_cache\_configs}
on real planner inputs captured from the observed run. The replay is bounded
and post-initialization. It is used to check whether the planner-stage
translation matches the runtime scheduler cache configuration. The current
evidence records:

\begin{itemize}
\item \texttt{planner\_matches\_runtime\_scheduler=true};
\item \texttt{runtime\_changed\_during\_replay=false}.
\end{itemize}

\subsection{Compute Advisory Metrics}

Cachepawl compares the observed reserved cache bytes with the useful byte
footprint implied by the translated hybrid cache plan. The diagnostic report
records observed reserved bytes, observed useful bytes, Cachepawl recommended
bytes for the advisory planner view, estimated advisory savings bytes,
overestimation ratio, and wasted fraction. These are planner-level advisory
metrics. They estimate over-reservation in the observed planner artifacts. They
are not serving-time VRAM measurements and do not imply throughput, latency,
accuracy, or quality effects.

\subsection{Package Artifact-Input Diagnosis}

\texttt{cachepawl diagnose-vllm} is the productized output of the phase. It
consumes an existing translated cache config and optional raw safe metadata,
then writes \texttt{report.json}, \texttt{summary.md}, and
\texttt{manifest.json}. The artifact-input CLI is intentionally lightweight. It
does not require vLLM, GPU access, CUDA, NVML, or model loading. It does not
modify vLLM, replace allocators, or enable runtime mutation.

\subsection{Inspect Mutation Contracts}

The final method step is not substitution. It is a contract audit that records
which runtime structures were safely observable and which were not. The current
evidence observed planner-stage replay match, request-to-block assignment,
worker tensor layout, and attention block-table/view metadata. Mamba
state-index and Mamba state tensor contracts remain blocked.

\section{Evaluation}

The evaluation has two evidence tracks. The Path C track measures
planner-level over-reservation in observed \texttt{vllm==0.21.0} cache-plan
artifacts for \texttt{Zyphra/Zamba2-2.7B-instruct}. The RTX 3060
planner-comparison track uses a deterministic synthetic planner harness for
\texttt{jamba-1.5-mini} workloads. Both tracks are advisory/planner-level only.

\subsection{Procedure and Non-Mutation Controls}

The Path C workflow observes a vanilla vLLM cache plan, translates it into the
Cachepawl schema, replays the vLLM planner on captured inputs, runs
\texttt{cachepawl diagnose-vllm}, and repeats the advisory diff over the
bounded 4-cell matrix. Planner-stage replay succeeded with
\texttt{planner\_matches\_runtime\_scheduler=true} and
\texttt{runtime\_changed\_during\_replay=false}.

All observations were read-only. The workflow did not modify vLLM, monkeypatch
private methods, replace allocators, return Cachepawl plans to vLLM, alter
scheduler behavior, or alter worker tensor layout. This validates the advisory
comparison against the observed planner/runtime config, but it does not
establish runtime substitution safety.

The planner-comparison pack is reproduced with:

\begin{verbatim}
UV_CACHE_DIR=/tmp/uv-cache uv run python \
  benchmarks/scripts/create_planner_comparison_pack.py \
  --output-dir benchmarks/results/rtx3060/planner-comparison \
  --seed 1 --num-requests 128 \
  --gpu-name "NVIDIA GeForce RTX 3060" \
  --gpu-total-bytes 12884901888
\end{verbatim}

The pack does not install vLLM, run serving, monkeypatch vLLM, replace
allocators, load a real model, run kernels, or measure latency and throughput.

\subsection{Metric Definitions}

\begin{itemize}
\item \texttt{reserved\_bytes}: bytes reserved by the observed or modeled
planner.
\item \texttt{useful\_bytes}: cache payload bytes before planner padding.
\item \texttt{estimated\_savings\_bytes = reserved\_bytes - useful\_bytes} for
Path C.
\item \texttt{overestimation\_ratio = reserved\_bytes / useful\_bytes}.
\item \texttt{wasted\_fraction = (reserved\_bytes - useful\_bytes) /
reserved\_bytes}.
\item \texttt{virtual\_oom}: whether a planner estimate exceeds the 12 GiB
target profile in the synthetic planner pack. It is not an observed runtime
OOM.
\end{itemize}

\subsection{Path C Advisory Matrix}

Table~\ref{tab:path-c-matrix} is advisory/diagnostic evidence only. It does not
report runtime mutation, throughput, serving, or VRAM improvement measurements.

\begin{table}[ht]
\centering
\small
\begin{tabular}{rrrrr}
\hline
\texttt{max\_model\_len} & \texttt{gpu\_memory\_utilization} & Reserved & Useful & Savings \\
\hline
2048 & 0.6 & 1,893,335,040 & 1,092,283,392 & 801,051,648 \\
2048 & 0.7 & 3,185,049,600 & 1,837,486,080 & 1,347,563,520 \\
4096 & 0.6 & 1,619,066,880 & 934,055,424 & 685,011,456 \\
4096 & 0.7 & 2,910,781,440 & 1,679,258,112 & 1,231,523,328 \\
\hline
\end{tabular}
\caption{Path C advisory matrix for \texttt{Zyphra/Zamba2-2.7B-instruct}. Values are bytes.}
\label{tab:path-c-matrix}
\end{table}

Across all four completed cells, the overestimation ratio stayed
1.7333734577189286 and the wasted fraction stayed 0.4230902777777778.
Estimated advisory savings ranged from 685,011,456 to 1,347,563,520 bytes. The
full matrix, including \texttt{num\_blocks}, cache group count, cache tensor
count, layer count, and status, is recorded in
\path{research/avmp/v2/evaluation/matrix_table.md}.

\subsection{RTX 3060 Planner-Only Comparison}

The planner-comparison pack reports deterministic planner estimates for three
synthetic hybrid KV/State workloads. The \texttt{vllm-style-padded} backend
models a uniform padded cache planner, not exact vLLM internals. The
\texttt{cachepawl-avmp} backend models a native KV-page plus SSM-block planner.
Table~\ref{tab:planner-comparison} is copied from
\path{benchmarks/results/rtx3060/planner-comparison/summary.md}. This is
planner evidence for the cache-shape claim, not runtime evidence for vLLM
serving behavior.

\begin{table}[ht]
\centering
\small
\begin{tabular}{llrrrrr}
\hline
Workload & Backend & Useful & Estimated & Over. & Waste & OOM \\
\hline
short-heavy & vllm-style-padded & 2,435,055,616 & 6,985,613,312 & 2.868770 & 0.651418 & false \\
short-heavy & cachepawl-avmp & 2,435,055,616 & 2,451,046,400 & 1.006567 & 0.006524 & false \\
long-heavy & vllm-style-padded & 41,967,550,464 & 165,116,116,992 & 3.934376 & 0.745830 & true \\
long-heavy & cachepawl-avmp & 41,967,550,464 & 41,983,672,320 & 1.000384 & 0.000384 & true \\
mixed & vllm-style-padded & 13,517,422,592 & 51,311,017,984 & 3.795917 & 0.736559 & true \\
mixed & cachepawl-avmp & 13,517,422,592 & 13,532,397,568 & 1.001108 & 0.001107 & true \\
\hline
\end{tabular}
\caption{RTX 3060 planner-only comparison. Useful and estimated columns are bytes.}
\label{tab:planner-comparison}
\end{table}

This supports a narrow planner claim: when the synthetic workload contains
heterogeneous KV and state-cache shapes, AVMP's planner model keeps estimated
bytes close to useful bytes, while the uniform padded planner model reserves
substantially more bytes. The result does not prove live serving memory
reduction or request-admission improvement. The pack does not run vLLM serving,
and \texttt{long-heavy} plus \texttt{mixed} remain virtual-OOM even under AVMP
because useful demand itself exceeds the 12 GiB target profile.

\subsection{Runtime Contracts and Claim Boundary}

Runtime contract observations resolved scheduler/KVCacheManager structure,
block usage metadata, worker tensor layout, live request-to-block assignment,
attention block-table/view metadata, and attention metadata builders. The Mamba
side remains gated: \texttt{mamba\_state\_idx} was reachable but empty for the
live request, no Mamba state tensors were safely reachable by stable runtime
attributes, and runtime cache config reported \texttt{mamba\_cache\_mode: none}.

These blockers are decisive for this phase. Without Mamba state-index and state
tensor contracts, the evidence supports observe/advisory reporting only. The
evidence does not show runtime allocator replacement, runtime cache
substitution, measured runtime VRAM reduction, throughput improvement, latency
improvement, quality improvement, or accuracy improvement.

\section{Limitations}

This report is scoped to a diagnostic/advisory systems artifact. Its main
limitations are:

\begin{itemize}
\item no runtime mutation: the evaluation does not replace vLLM allocators,
rewrite cache views, alter scheduler behavior, alter worker tensor layout, or
return Cachepawl plans to vLLM;
\item no serving performance or quality measurement: the evidence does not
include runtime VRAM reduction, throughput, latency, serving, quality, or
accuracy experiments;
\item narrow Path C scope: one observed model, vanilla \texttt{vllm==0.21.0},
\texttt{max\_model\_len} in \{2048, 4096\},
\texttt{gpu\_memory\_utilization} in \{0.6, 0.7\}, and
\texttt{max\_num\_seqs=1};
\item narrow planner-comparison scope: synthetic \texttt{jamba-1.5-mini}
workloads and a deterministic RTX 3060 target profile, not real model
inference;
\item local platform scope: the Path C artifacts are bounded to the recorded
local RTX 3060 12 GiB / WSL2 context and should not be generalized to
production serving clusters.
\end{itemize}

The Mamba-side mutation contracts remain blocked. The observed runtime cache
config reported \texttt{mamba\_cache\_mode: none}, so this run did not populate
the Mamba state-index/state-tensor paths needed for a view rewrite contract.
Because these contracts are blocked, the report cannot claim controlled
substitution readiness.

Future work must first improve Mamba state-index observability. A controlled
substitution experiment should only be considered after Mamba state paths are
observable through bounded read-only runtime paths, a supported vLLM integration
seam exists, or a default-off mutation probe with rollback controls and parity
checks is available. The next evaluation step is multi-model and
multi-workload advisory evaluation, still without conflating advisory savings
with measured runtime improvements.

\section{Conclusion}

Cachepawl Path C supports a narrow but useful systems claim: planner-level
hybrid cache over-reservation can be diagnosed from committed vLLM cache-plan
artifacts without modifying the serving stack. The current evidence includes a
bounded \texttt{Zyphra/Zamba2-2.7B-instruct} Path C matrix, a deterministic RTX
3060 planner-only comparison, and an artifact-input diagnostic CLI. Together
these support advisory reporting, not runtime replacement.

The technical ambition remains controlled substitution for hybrid cache
planning, but that is future work. Before this paper can claim live VRAM
reduction, throughput, latency, quality, accuracy, or allocator replacement, it
needs Mamba state-index and state tensor contracts, a default-off substitution
probe, and live serving measurements with rollback and parity controls.

\appendix
\section{Artifact and Reproducibility Checklist}

All paths are relative to the repository root. This checklist maps each report
claim to committed evidence and records the commands needed to verify the
venue-neutral LaTeX package.

\subsection{Claim Artifacts}

\begin{longtable}{p{0.42\linewidth}p{0.50\linewidth}}
\hline
Artifact & Claim supported \\
\hline
\path{research/avmp/v2/results/vllm-runtime-cache-diagnostic-cli/report.json} & User-facing advisory report with planner-level reserved/useful/savings metrics. \\
\path{research/avmp/v2/results/vllm-runtime-cache-diagnostic-cli/summary.md} & Human-readable advisory summary and non-mutation statement. \\
\path{research/avmp/v2/results/vllm-runtime-cache-diagnostic-cli/manifest.json} & Artifact-input mode, output list, and non-requirement flags for vLLM, CUDA, GPU, and NVML. \\
\path{research/avmp/v2/results/vllm-planner-stage-observation/translated_planner_stage_config.json} & Planner-stage replay output translated into Cachepawl's schema. \\
\path{research/avmp/v2/results/vllm-planner-stage-advisory-diff/diff_report.json} & Planner-stage advisory metrics and mutation-prevention field list. \\
\path{research/avmp/v2/results/vllm-planner-stage-advisory-diff/group_level_diff.json} & Per-cache-group advisory diff details. \\
\path{research/avmp/v2/evaluation/matrix_table.md} and \path{matrix_table.csv} & Four-cell advisory matrix and deterministic machine-readable metrics. \\
\path{benchmarks/results/rtx3060/planner-comparison/} & Deterministic RTX 3060 planner-only comparison pack. \\
\path{research/avmp/v2/evaluation/claim_summary.md} & Paper-facing claim boundary and unsupported claims. \\
\hline
\end{longtable}

\subsection{Reproducibility Checklist}

\begin{longtable}{p{0.27\linewidth}p{0.12\linewidth}p{0.53\linewidth}}
\hline
Item & Status & Evidence or command \\
\hline
Path C claim boundary documented & Ready & \path{research/avmp/v2/evaluation/claim_summary.md} \\
Path C matrix table committed & Ready & \path{research/avmp/v2/evaluation/matrix_table.md}, \path{research/avmp/v2/evaluation/matrix_table.csv} \\
Planner-comparison pack committed & Ready & \path{benchmarks/results/rtx3060/planner-comparison/} \\
Planner-comparison pack regeneration & Ready & Command listed below. \\
Diagnostic CLI tests & Ready & Command listed below. \\
Planner evidence tests & Ready & Command listed below. \\
Artifact path check & Ready & Confirm every path listed in this appendix exists in the repository. \\
Venue-specific formatting & TODO & Select target template, convert to the venue format, and add required author metadata and bibliography style. \\
\hline
\end{longtable}

Planner-comparison pack regeneration command:

\begin{verbatim}
UV_CACHE_DIR=/tmp/uv-cache uv run python \
  benchmarks/scripts/create_planner_comparison_pack.py \
  --output-dir benchmarks/results/rtx3060/planner-comparison \
  --seed 1 --num-requests 128 \
  --gpu-name "NVIDIA GeForce RTX 3060" \
  --gpu-total-bytes 12884901888
\end{verbatim}

Diagnostic CLI test command:

\begin{verbatim}
UV_CACHE_DIR=/tmp/uv-cache uv run pytest \
  tests/cli/test_diagnose_vllm.py -q
\end{verbatim}

Planner evidence test command:

\begin{verbatim}
UV_CACHE_DIR=/tmp/uv-cache uv run pytest \
  tests/bench/test_planner_comparison.py \
  tests/bench/test_vllm_path_c_advisory_matrix.py -q
\end{verbatim}

\subsection{Submission Readiness}

Ready: venue-neutral LaTeX draft with title, abstract, introduction, method,
evaluation, limitations, conclusion, artifact map, and reproducibility
checklist; compact evaluation tables for the Path C matrix and RTX 3060
planner-only comparison; explicit advisory/planner-level claim boundary
throughout the paper; committed evidence paths for numeric claims and product
artifacts.

Remaining before venue submission: choose the target venue format and page
limit; convert this simple article structure into the target template; add
venue-required author, anonymization, ethics, artifact, and bibliography
metadata; decide whether to keep the full artifact checklist in the main paper,
appendix, or supplemental material; run the venue-specific build once the
target template exists.

\bibliographystyle{plain}
\bibliography{references}

\end{document}
